\color{unidarkgreen}{\subsection{Блок отстройки от повреждений за трансформаторами ответвлений (ОПТО)}}
\color{black}

\begin{enumerate}

%%%%%%%%%%%%%%%%%%%%%%%%%%%%%%%%%%%%%%%%%%%%%%%%%% ОБШИЕ СВЕДЕНИЯ %%%%%%%%%%%%%%%%%%%%%%%%%%%%%%%%%%%%%%%%%%%%%%%%%%
\item Общие сведения

\begin{enumerate}

\item
Для предотвращения ложного срабатывания основной защиты при КЗ за трансформатором ответвления предусмотрен функциональный блок отстройки от повреждений за трансформаторами ответвлений (ОПТО).
В качестве измерительных органов ОПТО применяются три междуфазных прямонаправленных реле сопротивления и РНМНП-Р, блокируемое при выявлении броска намагничивающего тока.
При введенном программном ключе контроля защиты от ОПТО подготовка её цепи отключения разрешается только при срабатывании одного из реле сопротивления или ОНМНП-Р, что свидетельствует о наличии КЗ на линии, а не за трансформатором ответвления.

\item
Характеристика срабатывания РС задается уставками \textbf{«Rср РС-ОПТО»}, \textbf{«Хср РС-ОПТО»} и \textbf{«Ф1 РС-ОПТО»}.
Углы наклона характеристики срабатывания органа направленности, обеспечивающего направленность РС, задаются уставками \textbf{«Ф2 ОН-ОПТО»} и \textbf{«Ф3 ОН-ОПТО»}.
Для РС в составе ОПТО также предусмотрена возможность выреза в характеристике срабатывания для отстройки от нагрузки.
Характеристика срабатывания измерительного органа отстройки от нагрузки определяется уставками \textbf{«Rнг ОПТО»} и \textbf{«Фнг ОПТО»}.

\item
Характеристика срабатывания РС в составе ОПТО аналогична, представленной на рисунке \ref{pict_nvchz:RS_otkl_NVChZ}.

%\medskip
%
%\begin{figure}[H]
%\centering
%\includegraphics[width=0.4\textwidth,height=0.4\textheight,keepaspectratio]{\fbpath/01.01 ВЧ/561.1491 ОПТО/img/RS_OPTO.pdf}
%\captionof{figure}{Характеристика срабатывания РС в составе ОПТО}\label{pict_opto:RS_OPTO}
%\end{figure}
%
%\medskip

\item
Работа ОПТО может быть заблокирована при выявлении неисправности в цепях напряжения при введенном программном ключе \textbf{«БНН ОПТО»}.

\item
Вывод ОПТО из работы осуществляется программным переключателем \textbf{«ОПТО»}.
При введенном положении программного переключателя \textbf{«ОПТО»} вывод ОПТО из работы осуществляется при наличии сигнала \textit{«Вывод»} или \textit{«Вывод терминала»}. 

\item
Функциональная схема отстройки от повреждений за трансформаторами ответвлений представлена на рисунке \ref{pict_opto:OPTO}

\medskip

\begin{figure}[H]
\centering
\includegraphics[width=0.9\textwidth,height=0.9\textheight,keepaspectratio]{\fbpath/01.01 ВЧ/561.1491 ОПТО/img/OPTO.pdf}
\captionof{figure}{Функциональная схема отстройки от повреждений за трансформаторами ответвлений}
\label{pict_opto:OPTO}
\end{figure}

\medskip

\item 
Уставочные параметры функции ОПТО представлены в таблице \ref{app_set:opto}.

\end{enumerate}

\end{enumerate}
